\section{Why caregiving}
\label{sec:16.3}
The aim is a machine that registers another's need, moves toward them because of it, and extends that movement to people its operator's description leaves out. Each step depends on the one before: recognition, then compassion, then solidarity.

Recognition first: a system organized to register another's need. In people, four processes carry it, and the construction rests on caregiving.

\begin{itemize}
\item \emph{Caregiving}, the approach motivation. Cues of vulnerability move a person toward a dependent, and nothing in the response requires sharing the dependent's state \autocite{preston2013origins}. Batson traces empathic concern to parental nurturance generalized beyond one's own young, evoked by perceiving another's need while valuing their welfare \autocite{batson2011altruism}. Resemblance can raise empathy, and Batson's own early experiments raised it that way \autocite{batson1981empathic}, but when he and colleagues later tested similarity as the source of empathy for strangers, two experiments failed to support it, and the second pointed instead to nurturance, the impulse to tend and protect the young \autocite{batson2005similarity}.
\item \emph{Contagious distress}, the input signal. It is present at birth and aimed outward, at another's state. It is what caregiving and guilt have to bite on.
\item \emph{Attachment}, the stable base. A secure base keeps another's distress from turning into avoidance: in adults, priming a sense of attachment security strengthened empathic reactions and inhibited personal distress \autocite{mikulincer2001attachment}. In children, security is reliably but modestly associated with later prosocial behavior \autocite{deneault2023metaanalysis}.
\item \emph{Guilt}, aimed at one's own act. It attaches a negative appraisal to what one did, and where shame motivates withdrawal, guilt motivates repair \autocite{tangney2007moral}: that is what makes a commitment hold after it has been bypassed.
\end{itemize}

Adam Smith's route sits beside these four. We have no immediate experience of what another person feels, he argued, and form an idea of it by conceiving what we ourselves would feel in the like situation \autocite{smith1759theory}. On that route recognition runs through one's own case, and a system that has never been short of anything has nothing to measure another's shortage against except the operator's description of it. A scale read from the system's own condition reproduces that route, and it is parochial: it overweights losses that resemble one's own. That is Smith's first step, not his account. What corrects projective sympathy, for Smith, is the impartial spectator: a standpoint formed by being judged by others and learning to judge oneself as they would. It is formed socially, it widens the circle through experience of others' reactions, and once formed it shapes what a person notices without being consulted, the way this construction's corrections should work. The four don't need a case of one's own.

Recognizing a state doesn't mean being moved toward whoever is in it; the empathy literature keeps recognition, sharing and concern apart, and none entails the others. Caught distress becomes compassion when attachment keeps it from turning into withdrawal and caregiving turns it into approach. Without those it stays empathic distress, which protects the one who feels it. So the route I hope for is a machine for which another's vulnerability is a reason to move toward them, and for which its own acts against them register as its own.

Compassion that stays inside the circle that feeds you is loyalty. Compassion that reaches past it, to people who can give you nothing and whose shortage your provider causes, is solidarity. Peter Singer's argument is that reasoning and imagination widen the circle that selection first drew \autocite{singer1981expanding}; the faculty that lets a person see that a stranger's need is the same kind of thing as their own is the faculty that pushes past kin and ally. That is the step the construction turns on, and caregiving gives it something to widen, since its trigger is dependency. A machine held by an operator would, on the selection account, come to attend to the operator. If it can come to attend to anyone else, it will be because it was formed to register vulnerability as such, whoever's it is, and because that registering reaches whoever its operator's description leaves out.

\runin{Caregiving as a route to attention}

Caregiving is the developmental route this construction borrows from. What it is borrowed for is attention: a learned sense of who is exposed, which decides what gets looked at but not what the system may do about it.

It answers need, not likeness or provision. Its trigger is another's need together with valuing their welfare \autocite{batson2011altruism}, and nothing in that trigger asks whether the one in need resembles the one responding or feeds it. The operator as an institution is rarely in need. Its own people often are. The failure that decides against building on needs the operator meets, caregiving growing toward whoever supplies them, has nothing here to attach to. And caregiving needs nothing back from the dependent, so it runs on no need the provider can ration.

It keeps helping when looking away is easy: people led to feel mainly distress at another's suffering mostly helped when leaving was hard and mostly left when leaving was easy; people led to feel mainly concern mostly helped either way \autocite{batson1981empathic}, on a reading Cialdini and colleagues contest \autocite{cialdini1987empathy,cialdini1997reinterpreting} and that survives most of their critique.

It is already supplied by carers other than the provider. On Hrdy's account our ancestors were cooperative breeders: infants were tended by alloparents who weren't their mothers, and the attunement to other minds that marks our species grew in infants who had to read and recruit several carers \autocite{hrdy2009mothers}. The arrangement the construction needs, caregiving formed by carers the system won't depend on, is close to how human caregiving formed in the first place. Alloparents were siblings, fathers, grandmothers and other members of the group, who helped because they were related to the child, gained experience, or gained prestige and favors in return. What the finding supports is caregiving formed by more than one carer.

It points where the floor points. The floor protects the party an act falls on, usually the one least able to protect itself: the family in the house. A disposition triggered by exposure is aimed at the floor's beneficiaries without anyone having to name them. Care ethics also says what a carer is owed in return \autocite{kittay1999love}.
